\section{Cinco lecciones prácticas de cargar por adelantado}

Las siguientes cinco imágenes no demuestran repetidamente lo mismo, sino que responden secuencialmente a cinco objeciones: ¿Es efectivo inmediatamente después del preentrenamiento? ¿El SFT elimina los beneficios? ¿Sobrefitentan los datos de alta calidad? ¿Puede compensarse con más cómputo de SFT o más datos tardíos? ¿Cuánto queda después de completar RLVR? Cada gráfico debe verificar primero la comparación, luego leer la altura de las barras.

\subsection{Lección 1: Beneficios inmediatos tras el preentrenamiento}

La primera lección revisa el punto temporal más directo: antes de entrar en SFT o RL, ¿añadir datos de razonamiento altera ya al modelo base? Esta verificación elimina todos los factores posteriores al entrenamiento y observa si la asignación de datos durante el preentrenamiento es suficiente para generar una diferencia medible. Se comparan las bases con y sin razonamiento tras completar el preentrenamiento; la mejora relativa promedio en Ampere es de aproximadamente 16\%. Esto indica que los datos de razonamiento no solo enseñan al modelo a producir un formato específico, sino que pueden cambiar la representabilidad de la base para tareas matemáticas, científicas y de razonamiento general. Aún no prueba que el modelo haya generado internamente un proceso de razonamiento fielmente interpretable.

\lecturefigure{slide-017.jpg}{Beneficios inmediatos de añadir datos de razonamiento en la fase de preentrenamiento}{Grabación oficial de Stanford, 00:15:48--00:16:36}

\begin{knowledgebox}{Lectura de la figura: orden de comparación}
Verificar primero que ambos grupos de modelos son consistentes en arquitectura, presupuesto de tokens y mezcla de preentrenamiento estándar; luego observar las barras azul y verde en cada benchmark; finalmente revisar la mejora relativa promedio. Si un punto de partida es muy bajo, pequeños cambios absolutos pueden generar porcentajes relativos grandes, por lo que se debe observar tanto la puntuación cruda como la ganancia relativa.
\end{knowledgebox}

\subsection{Lección 2: El SFT no elimina las ventajas}

Que haya beneficios inmediatos no es suficiente, porque los modelos de producción generalmente continúan recibiendo SFT; por ello, la segunda lección aplica el mismo entrenamiento posterior a ambas fundaciones para verificar si la ventaja es solo un fenómeno efímero. Tras someter a dos modelos base idénticos al mismo SFT, la base con razonamiento mantiene una ventaja relativa promedio de aproximadamente 9.3\%. La conclusión pedagógica aquí es "la fundación importa": el entrenamiento posterior no reinicia completamente el modelo, sino que sigue optimizando sobre su representación base. Las ganancias tempranas pueden disminuirse, mantenerse o amplificarse en tareas específicas; se debe verificar por tarea.

\lecturefigure{slide-018.jpg}{Las ganancias iniciales de razonamiento persisten tras el SFT}{Grabación oficial de Stanford, 00:16:36--00:17:14}

\teachervoice{De 00:16:36 a 00:17:14, la oradora usa repetidamente la expresión "washed away". Le interesa no solo la puntuación base, sino si esa misma ventaja sigue siendo observable tras el entrenamiento posterior; este es el criterio clave para determinar si cargar por adelantado vale la pena desde una perspectiva de ingeniería.}

\subsection{Lección 3: Los datos de alta calidad pueden tener un efecto latente}

La tercera lección aborda un problema más sutil: si los datos de alta calidad no elevan inmediatamente las puntuaciones en la evaluación base, ¿significa que no tienen valor? El curso clasifica los datos de razonamiento en SHQ, LDQ y LMQ. SHQ es pequeño pero de alta calidad; LDQ es grande y diverso pero de menor calidad; LMQ combina ambos. Al finalizar el preentrenamiento, LMQ y LDQ pueden ser casi idénticos; tras el SFT, LMQ muestra una mejora adicional de aproximadamente 4.25\%. La señal de alta calidad no se refleja inmediatamente en los benchmarks base, sino que puede ofrecer mejor plasticidad para optimizaciones futuras.

\lecturefigure{slide-019.jpg}{El beneficio potencial de los datos de preentrenamiento de alta calidad se activa tras el SFT}{Grabación oficial de Stanford, 00:17:14--00:18:48}

\begin{knowledgebox}{Glosario: SHQ, LDQ, LMQ}
\begin{tabularx}{\textwidth}{L{0.16\textwidth}L{0.23\textwidth}X}
\toprule
Abreviatura & Atributo de los datos & Función pedagógica \\
\midrule
SHQ & Pequeño, Alta Calidad & Verificar si una cantidad pequeña de datos filtrados con precisión proporciona estructura adicional \\
LDQ & Grande, Diverso, Baja Calidad & Verificar la base de escala y diversidad \\
LMQ & Combinación de SHQ + LDQ & Verificar si el componente de alta calidad muestra un efecto latente tras el entrenamiento posterior \\
\bottomrule
\end{tabularx}
\end{knowledgebox}

\subsection{Lección 4: Duplicar cómputo tardío o desplazar datos no logra alcanzarlo}

En la figura izquierda se le dan más épocas de SFT a la base sin razonamiento; en la derecha se fija el presupuesto total de datos de razonamiento y se compara "todo en SFT" frente a "preentrenamiento + asignación en SFT". En ambos casos, el modelo cargado por adelantado mantiene su ventaja. Esto refuta respectivamente las explicaciones de que "solo falta más entrenamiento posterior" o que "basta con usar los mismos datos más tarde".

\lecturefigure{slide-020.jpg}{Efecto fuerte de la fundación bajo condiciones de cómputo y datos emparejados}{Grabación oficial de Stanford, 00:18:48--00:20:54}

\begin{importantbox}{Por qué realizar ambos tipos de emparejamiento simultáneamente}
El emparejamiento de cómputo controla la cantidad de entrenamiento computacional; el emparejamiento de datos controla el presupuesto de muestras; ambos resuelven factores de confusión distintos. Si la condición temprana usa más datos y consume más FLOPs, no se puede atribuir el beneficio únicamente al momento. Un reporte riguroso debe listar al menos los tokens, épocas, pasos del optimizador, longitud de secuencia y el costo adicional de rollout/SFT.
\end{importantbox}

\subsection{Lección 5: Persiste una ventaja duradera tras RLVR}

Que siga liderando después del SFT no garantiza que RLVR no vuelva a nivelar ambas trayectorias; la quinta lección verifica entonces el pipeline final entregable, en lugar de detenerse en un checkpoint intermedio. La ventaja relativa promedio de Ampere frente a Volta es de aproximadamente 19\%, con diferencias relativas aún mayores en tareas matemáticas difíciles como AIME. El gráfico apoya la idea de que "el razonamiento temprano puede generar una ventaja duradera y acumulativa", pero no indica que todas las tareas crezcan al mismo ritmo, ni que el entrenamiento posterior deje de ser importante; es precisamente el entrenamiento posterior idéntico lo que expone las diferencias entre fundaciones.

\lecturefigure{slide-021.jpg}{Cargar por adelantado mantiene la ventaja tras el entrenamiento completo posterior}{Grabación oficial de Stanford, 00:20:54--00:21:48}

\begin{warningbox}{Una mejora relativa grande no implica cercanía a la puntuación máxima}
Cuando las tareas como AIME tienen una línea base baja, varias puntos de mejora absoluta pueden corresponder a una ganancia relativa muy grande. Las decisiones de ingeniería deben reportar simultáneamente la puntuación absoluta, el promedio, el número de ensayos y la varianza; no basta con recortar el porcentaje relativo más alto para afirmar una mejora general de capacidades.
\end{warningbox}

\subsection{¿Qué detalles de ingeniería fuera de los artículos completó la sesión de preguntas y respuestas?}

La sesión de preguntas intermedias aclaró que el uso de datos de razonamiento se basa en conjuntos de datos existentes en la comunidad y definiciones de trabajo comunes, donde las muestras típicas son pregunta, rastro de razonamiento largo y respuesta final, con un sesgo experimental hacia STEM. La calidad puede operacionalizarse mediante intensidad de filtrado, complejidad, longitud o etiqueta de dificultad; la optimización de la receta es un proceso híbrido de estimación automática y ensamblaje manual. La oradora también especificó que invertir el orden de las fases resultó peor en sus estudios de abstracción.

\teachervoice{De 00:22:00 a 00:27:58, la oradora no empaquetó "razonamiento" como un concepto unificado transversal. Reconoció que los datos usados se inclinan hacia matemáticas, código y STEM, y el blend final tampoco es totalmente automático. Estas limitaciones determinan los límites de extrapolación de los resultados.}

\begin{knowledgebox}{Práctica profesional: cómo registrar una receta de datos reproducible}
Se debe guardar al menos la versión de la fuente de datos, las claves de deduplicación, filtros y umbrales, rúbrica de calidad, mapeo de dominios, número de tokens únicos, épocas objetivo, pesos por fase, límite de fases, semilla aleatoria y configuración de evaluación proxy. Publicar solo el porcentaje final no permite distinguir diferencias causadas por estimaciones de calidad, implementación de muestreo o orden.
\end{knowledgebox}

\subsection{Resumen de la sección}

Las cinco lecciones establecen la cadena de evidencia para cargar por adelantado: los beneficios aparecen inmediatamente, no desaparecen tras el SFT, las señales de alta calidad pueden manifestarse con retraso, más cómputo o datos tardíos no logran igualar completamente y la ventaja persiste incluso después del RLVR completo. La siguiente parte profundiza aún más: además de proporcionar datos de razonamiento antes, ¿podemos hacer que el modelo explore activamente sus propios pensamientos durante el preentrenamiento?
